\documentclass[12pt]{article}
\usepackage{notes-project}
\usepackage{braket}

\begin{document}

\begin{center}
  {\LARGE\bfseries Chapter 4: Quantum Mechanics in 3D}\\[0.25em]
  {\large Quantum Mechanics}\\[0.15em]
\end{center}

% Key definitions from Griffiths Ch. 4 typed as needed for cross-references.

\section{Angular Momentum}

\begin{dfn}[Angular Momentum Operator]{angular-momentum}
  The \emph{angular momentum operator} is
  \[
    \hat{\vec{L}} \;:=\; \hat{\vec{r}} \times \hat{\vec{p}},
  \]
  with components
  \[
    \hat{L}_x = \hat{y}\hat{p}_z - \hat{z}\hat{p}_y, \quad
    \hat{L}_y = \hat{z}\hat{p}_x - \hat{x}\hat{p}_z, \quad
    \hat{L}_z = \hat{x}\hat{p}_y - \hat{y}\hat{p}_x.
  \]
\end{dfn}

\section{Atomic Eigenstates}

\begin{dfn}[Atomic Eigenstate]{atomic-eigenstate}
  An \emph{atomic eigenstate} $\ket{\psi_{n\ell m}}$ (also written $\ket{n,\ell,m}$) is the simultaneous eigenstate of $\hat{H}$, $\hat{L}^2$, and $\hat{L}_z$ with eigenvalues $E_n$, $\ell(\ell+1)\hbar^2$, and $m\hbar$ respectively. Its position-space wavefunction separates as
  \[
    \psi_{n\ell m}(r, \theta, \phi) = R_{n\ell}(r)\, Y_\ell^m(\theta, \phi),
  \]
  where $R_{n\ell}$ is the radial function and $Y_\ell^m$ is the spherical harmonic.
\end{dfn}

\begin{thm}[Parity of Spherical Harmonics]{spherical-harmonic-parity}
  Under $(\theta, \phi) \to (\pi - \theta,\, \phi + \pi)$,
  \[
    Y_\ell^m(\pi - \theta,\, \phi + \pi) = (-1)^\ell\, Y_\ell^m(\theta, \phi).
  \]
\end{thm}

\begin{proof}
  {\color{blue}\textbf{TODO}}
\end{proof}

\end{document}
